\section{核心理论与模型架构}
\label{sec:theory}

本章介绍实验中用到的五种模型，涵盖统计机器学习、深度序列建模、预训练编码器和大语言模型四个阶段。各模型在特征提取方式、上下文建模能力和参数更新策略上各有不同。

\subsection{基于统计特征的 SVM 分类模型}
在深度学习流行之前，支持向量机（Support Vector Machine, SVM）一直是文本分类的标准方案。文本通过词袋模型（Bag-of-Words）转为高维稀疏向量后，送入 SVM 寻找最大间隔超平面进行分类。

为抑制高频停用词的影响，本实验用 TF-IDF（Term Frequency-Inverse Document Frequency）对文本做特征加权。词 $t$ 在文档 $d$ 中的 TF-IDF 权重为：
\begin{equation}
	\text{TF-IDF}(t, d) = \text{TF}(t, d) \times \log \left( \frac{N}{\text{DF}(t) + 1} \right)
\end{equation}
其中，$\text{TF}(t, d)$ 是词频，$N$ 是文档总数，$\text{DF}(t)$ 是包含该词的文档数。这个公式会压低"的""了"等高频词的权重，抬高有情感区分度的特征词。

SVM 在这个 TF-IDF 空间中求解最大间隔超平面。引入软间隔和惩罚参数 $C$ 后，对偶优化问题为：
\begin{equation}
	\begin{aligned}
		\min_{\alpha} \quad & \frac{1}{2} \sum_{i=1}^{m} \sum_{j=1}^{m} \alpha_i \alpha_j y_i y_j K(x_i, x_j) - \sum_{i=1}^{m} \alpha_i \\
		\text{s.t.} \quad & \sum_{i=1}^{m} \alpha_i y_i = 0, \quad 0 \leq \alpha_i \leq C
	\end{aligned}
\end{equation}
$\alpha_i$ 是拉格朗日乘子，$K(x_i, x_j)$ 是核函数。文本特征通常近似线性可分，用线性核的 SVM 在特征数远大于样本数时泛化能力很好，且优化目标为凸二次规划，收敛速度快。

以评论"酒店环境很好，服务也不错"为例，分词后得到[酒店, 环境, 很, 好, 服务, 也, 不错]。SVM 为每个词分配一个权重：情感词"好""不错"获得高正向权重，"酒店""环境""服务"等中性词权重接近零，"很""也"等停用词因 IDF 压减而被忽略。最终决策值由所有词的加权和决定。如果在这条评论中插入否定词——变为"酒店环境不是很好，服务也不怎么样"——SVM 无法从词袋中区分"好"在两个上下文中的不同语义，因为它的输入只是词频向量，词序信息被完全丢弃。

\subsection{基于深度序列建模的 Bi-LSTM 网络}
词袋模型丢掉了词序。但中文情感经常依赖前后文的逻辑关系——比如双重否定、转折词。长短期记忆网络（LSTM）通过门控机制缓解了 RNN 的梯度消失和爆炸问题。

LSTM 单元包含遗忘门 $f_t$、输入门 $i_t$ 和输出门 $o_t$。在时间步 $t$，输入词嵌入向量 $x_t$，细胞状态 $C_t$ 和隐藏状态 $h_t$ 的更新为：
\begin{align}
	f_t &= \sigma(W_f \cdot [h_{t-1}, x_t] + b_f) \label{eq:lstm_forget} \\
	i_t &= \sigma(W_i \cdot [h_{t-1}, x_t] + b_i) \label{eq:lstm_input} \\
	\tilde{C}_t &= \tanh(W_C \cdot [h_{t-1}, x_t] + b_C) \label{eq:lstm_candidate} \\
	C_t &= f_t \odot C_{t-1} + i_t \odot \tilde{C}_t \label{eq:lstm_cell} \\
	o_t &= \sigma(W_o \cdot [h_{t-1}, x_t] + b_o) \label{eq:lstm_output} \\
	h_t &= o_t \odot \tanh(C_t) \label{eq:lstm_hidden}
\end{align}

单向 LSTM 只看到一个方向的信息。中文评论的转折词常出现在句中甚至句末，只靠前向编码会丢失后文信息。本实验用双向 LSTM（Bi-LSTM），把前向隐藏状态 $\overrightarrow{h_t}$ 和后向隐藏状态 $\overleftarrow{h_t}$ 拼起来：
\begin{equation}
	h_t = [\overrightarrow{h_t} \oplus \overleftarrow{h_t}]
\end{equation}
双向结构使每个词都能同时感知左右两侧的语境，有助于捕捉中文句法和情感翻转。

以"虽然服务慢但味道真不错"为例，单向 LSTM 在读到"不错"时只能看到"虽然服务慢但味道真不错"的前文，但不知道句末是否还有反转。Bi-LSTM 同时从后向前编码：后向 LSTM 先读到"不错"，再回溯"味道真"和"但"，将"但"识别为转折信号，从而确认"不错"是最终的情感基调。在处理"不是不好"这样的双重否定时，Bi-LSTM 的前向通道捕捉到"不是"，后向通道回溯到"不好"，两个方向的隐状态拼接后包含了完整的否定链信息，分类层据此做出正向判断。

\subsection{预训练编码器 BERT 模型}
Transformer 不再依赖循环时序，而是通过自注意力（Self-Attention）实现并行计算和全局上下文建模。BERT（Bidirectional Encoder Representations from Transformers）基于多层双向 Transformer 编码器构建，其核心操作为缩放点积自注意力：

\begin{equation}
	\text{Attention}(Q, K, V) = \text{softmax}\left(\frac{Q K^T}{\sqrt{d_k}}\right) V
\end{equation}

这部分数学原理在第~\ref{sec:mhsa_mechanism}~节展开。BERT 通过掩码语言模型（Masked Language Model, MLM）在中文语料上双向预训练，获得了深层上下文语义表征。微调时，在输入开头放一个 \texttt{[CLS]} 标记，经多层自注意力交互后，取该标记的最终隐状态 $C \in \mathbb{R}^H$ 作为整句语义表示，送入线性分类层，用交叉熵损失优化。

BERT 处理中文的优势在于其预训练语料覆盖了海量真实文本。以"这家店的蛋糕甜而不腻"为例：Bi-LSTM 需要从标注数据中学到"甜而不腻"整体为正面的模式，而 BERT 已在预训练中见过大量类似表达，知道"甜而不腻"在中文语境中是称赞。当遇到训练集中未见的表达——例如"入口即化"——BERT 通过 MLM 预训练获得的语义知识仍能为 \texttt{[CLS]} 标记生成合理表征，而 Bi-LSTM 可能将其视为未知组合，只能靠训练集中的有限模式外推。

\subsection{Qwen 大规模语言模型与注意力机制}
\label{sec:mhsa_mechanism}

本实验还引入了通义千问（Qwen）。Qwen 是阿里云发布的开源大模型系列，基于 Transformer 解码器（Decoder-only）架构，采用因果自注意力（Causal Self-Attention）逐 Token 做自回归预测。实验选用 Qwen-7B：70 亿参数、32 层 Transformer、32 注意力头、隐藏维度 4096，规模远大于 BERT-base（约 1.1 亿参数）。

\subsubsection{缩放点积自注意力 (Scaled Dot-Product Attention)}
自注意力的本质是计算序列中每对 Token 之间的相关性。输入向量序列通过三个可学习权重矩阵映射为查询向量 $Q$、键向量 $K$、值向量 $V$：
\begin{equation}
	Q = XW^Q, \quad K = XW^K, \quad V = XW^V
\end{equation}
其中，$W^Q, W^K, W^V \in \mathbb{R}^{d_{\text{model}} \times d_k}$。用缩放点积计算 $Q$ 和 $K$ 的相似度：
\begin{equation}
	\text{Attention}(Q, K, V) = \text{softmax}\left(\frac{Q K^T}{\sqrt{d_k}}\right) V
\end{equation}
除以 $\frac{1}{\sqrt{d_k}}$ 是为了防止点积值过大，导致 Softmax 进入梯度平坦区。

\subsubsection{多头机制的并行表征 (Multi-Head Mechanism)}
单头注意力只在一个子空间内建模。多头机制把 $Q, K, V$ 投影到 $h$ 个不同的低维子空间，并行计算。比如在做情感分析时，一个头可能关注否定词，另一个关注转折连词。

对第 $i$ 个头：
\begin{equation}
	\text{head}_i = \text{Attention}(Q W_i^Q, K W_i^K, V W_i^V)
\end{equation}
所有头的输出拼接后经过 $W^O$ 做线性变换：
\begin{equation}
	\text{MultiHead}(Q, K, V) = \text{Concat}(\text{head}_1, \dots, \text{head}_h) W^O
\end{equation}

\subsubsection{情感分析中的全局特征聚合}
在情感分类中，全连接的注意力图消除了序列距离的限制。比如输入"虽然环境很差，但味道真的一流"，MHSA 让句末的"一流"能直接和句首"但"转折后的语义建立关联，从而定位正向信号。这是 RNN 和浅层模型做不到的全局依赖建模。

\subsection{参数高效微调技术 LoRA}
Qwen-7B 虽然在零样本下已具备一定的判断力，但对 70 亿参数做全量微调的显存和计算开销远超出常规实验条件。

LoRA（Low-Rank Adaptation）的思路是：大模型适配下游任务时，权重更新量 $\Delta W$ 具有低秩性质。因此可以在冻结主干权重 $W_0 \in \mathbb{R}^{d \times k}$ 的同时，旁路插入低秩分解矩阵：
\begin{equation}
	h = W_0 x + \Delta W x = W_0 x + B A x
\end{equation}
$B \in \mathbb{R}^{d \times r}$ 和 $A \in \mathbb{R}^{r \times k}$ 是可训练的低秩矩阵，$r \ll \min(d, k)$。反向传播只更新 $A$ 和 $B$ 的梯度。本实验取 $r=16$，可训练参数压缩到全量微调的千分之几。在单张 GPU 上即可完成 7B 级别模型的情感分类适配，推理时做权重合并，延迟与原模型一致。
